% GENERATED FILE. Edit canonical lessons, then run npm run generate:books.
% canonical-source-hash: fnv1a64:ea7ece39c62617f9
% canonical-lessons: SA-C267-school-recall, SA-C267-exam-time, SA-C267-homework-plan, SA-C267-complete-school-plan, SA-C267-mixed-school-practice, SA-C267-school-plan-retrieval

\chapter{One small school-day plan}
\label{ch:sa-school-day-plan}
\hlchaptermodality{267}

\begin{chapteropening}
\textbf{By the end of this chapter:} \emph{I can ask when an exam will be and state a short homework and study plan.}

Complete chapter 267 by independently asking and answering an exam-time question and stating two schoolwork plans.
\end{chapteropening}

\section[(school-word recall)]{The school words are already yours}

\label{lesson:SA-C267-school-recall}



\begin{quote}
Say “a game,” then state the cue carried by the complete first-person future forms, then say “I buy.” (\emph{Khelā; the taught first-person-singular future cue; krīṇāmi.})
\end{quote}

\subsection*{Your turn}
\begin{itemize}
\raggedright
  \item \emph{Recall:} \textbf{\sk{विद्यालयः}; \sk{छात्रः}; \sk{अध्यापकः}} — \emph{vidyālayaḥ; chātraḥ; adhyāpakaḥ} — school; student; teacher
  \item \emph{Recall:} \textbf{\sk{पाठः}; \sk{परीक्षा}} — \emph{pāṭhaḥ; parīkṣā} — lesson; exam
  \item \emph{Recall:} \textbf{\sk{समयसारणी}; \sk{कक्षा}} — \emph{samayasāraṇī; kakṣā} — timetable; classroom
  \item \emph{Recall:} \textbf{\sk{अंकः}} — \emph{aṅkaḥ} — a mark or number
  \item \emph{Recall:} \textbf{\sk{गृहकार्यम्}; \sk{उत्तरम्}} — \emph{gṛhakāryam; uttaram} — homework; answer
\end{itemize}

\begin{grammarlens}[title={What you've built}]
The institution, its people, the timetable, the classroom, and the assessment words were already in the book. This chapter makes them carry one small plan.
\end{grammarlens}

\begin{tcolorbox}[breakable,colback=teal!4,colframe=teal!35!black,title={Before you move on}]
Say five pieces without looking: school, exam, timetable, classroom, and mark.
\end{tcolorbox}
\section[(build one bounded school plan)]{Build one bounded school plan}

\label{lesson:SA-C267-exam-time}



\begin{quote}
Ask the familiar market question and give its short rupee answer. (\emph{Mūlyaṃ kiyat? Viṃśatiḥ rūpyakāṇi.})
\end{quote}

\begin{grammarlens}[title={Grammar Lens: four complete turns}]
Put four known pieces together:

\begin{quote}
\textbf{\sk{परीक्षा} \sk{कदा} \sk{भविष्यति}?}
\end{quote}

\begin{quote}
\emph{Parīkṣā kadā bhaviṣyati?}
\end{quote}

\begin{quote}
When will the exam be?
\end{quote}

Answer from the smallest possible timetable:

\begin{quote}
\textbf{\sk{श्वः} \sk{परीक्षा} \sk{भविष्यति।}}
\end{quote}

\begin{quote}
\emph{Śvaḥ parīkṣā bhaviṣyati.}
\end{quote}

\begin{quote}
The exam will be tomorrow.
\end{quote}

This is one useful school-time pattern. It does not teach a new future paradigm or a complete timetable system. Complete the bounded plan with two whole schoolwork frames:

\begin{quote}
\textbf{\sk{गृहकार्यं} \sk{करिष्यामि।}}
\end{quote}

\begin{quote}
\emph{Gṛhakāryaṃ kariṣyāmi.}
\end{quote}

\begin{quote}
I will do the homework.
\end{quote}

\begin{quote}
\textbf{\sk{पाठं} \sk{पठिष्यामि।}}
\end{quote}

\begin{quote}
\emph{Pāṭhaṃ paṭhiṣyāmi.}
\end{quote}

\begin{quote}
I will read or study the lesson.
\end{quote}

Own those two lines as complete frames. In particular, \textbf{\sk{पाठं}} belongs to the whole second frame here; do not use it to infer a masculine accusative rule.
\end{grammarlens}

\subsection*{Your turn}
Say these aloud:

\begin{itemize}
\raggedright
  \item ask when the exam will be
  \item answer that it will be tomorrow
  \item the complete homework frame
  \item the complete study frame
\end{itemize}

\begin{tcolorbox}[breakable,colback=teal!4,colframe=teal!35!black,title={Before you move on}]
Ask and answer the exam-time turn, then say both complete schoolwork frames from memory.
\end{tcolorbox}
\section[(one homework and study plan)]{Retrieve your homework and study plan}

\label{lesson:SA-C267-homework-plan}



\begin{quote}
Say “a lesson,” then “homework.” (\emph{Pāṭhaḥ; gṛhakāryam.})
\end{quote}

\begin{grammarlens}[title={Grammar Lens: two complete owned frames}]
Retrieve the two whole frames introduced in the preceding lesson:

\begin{quote}
\textbf{\sk{गृहकार्यं} \sk{करिष्यामि।}}
\end{quote}

\begin{quote}
\emph{Gṛhakāryaṃ kariṣyāmi.}
\end{quote}

\begin{quote}
I will do the homework.
\end{quote}

\begin{quote}
\textbf{\sk{पाठं} \sk{पठिष्यामि।}}
\end{quote}

\begin{quote}
\emph{Pāṭhaṃ paṭhiṣyāmi.}
\end{quote}

\begin{quote}
I will read or study the lesson.
\end{quote}

The school-plan atom owns both complete lines. This lesson does not derive \textbf{\sk{पाठं}} from \textbf{\sk{पाठः}}, generalise a masculine accusative rule, or add a new future form.
\end{grammarlens}

\subsection*{Your turn}
Say these aloud:

\begin{itemize}
\raggedright
  \item “I will do the homework”
  \item “I will study the lesson”
  \item both plans after the exam-time answer
\end{itemize}

\begin{tcolorbox}[breakable,colback=teal!4,colframe=teal!35!black,title={Before you move on}]
State the homework plan and the study plan without looking.
\end{tcolorbox}
\section[(complete one school-day plan)]{Join the pieces into one school-day plan}

\label{lesson:SA-C267-complete-school-plan}



\begin{quote}
Say “a fort,” then “I will do.” (\emph{Durgam; kariṣyāmi.})
\end{quote}

\subsection*{Your turn}
First name the setting and its useful records: \textbf{\sk{विद्यालयः}}, school; \textbf{\sk{कक्षा}}, classroom; \textbf{\sk{समयसारणी}}, timetable; \textbf{\sk{अंकः}}, mark; \textbf{\sk{उत्तरम्}}, answer. Name the people: \textbf{\sk{छात्रः}}, student; \textbf{\sk{अध्यापकः}}, teacher.

Then say the working plan:

\begin{enumerate}
\raggedright
  \item \textbf{\sk{परीक्षा} \sk{कदा} \sk{भविष्यति}?} — When will the exam be?
  \item \textbf{\sk{श्वः} \sk{परीक्षा} \sk{भविष्यति।}} — The exam will be tomorrow.
  \item \textbf{\sk{गृहकार्यं} \sk{करिष्यामि।}} — I will do the homework.
  \item \textbf{\sk{पाठं} \sk{पठिष्यामि।}} — I will study the lesson.
\end{enumerate}

\begin{grammarlens}[title={What you've built}]
The nouns were already yours. Now a timetable question, answer, and two plans hold them together without adding a new paradigm.
\end{grammarlens}

\begin{tcolorbox}[breakable,colback=teal!4,colframe=teal!35!black,title={Before you move on}]
Say the question, the answer, the homework plan, and the study plan.
\end{tcolorbox}
\section[(independent school plan, first round)]{Independent school plan: round one}

\label{lesson:SA-C267-mixed-school-practice}



\begin{quote}
Say “a palace,” then “I will read or study.” (\emph{Prāsādaḥ; paṭhiṣyāmi.})
\end{quote}

\subsection*{Your turn: answer before the check}
No Sanskrit model is shown before you answer.

You are a student reading a school timetable. Ask when the exam will be. Answer that it will be tomorrow. State that you will do the homework and study the lesson. Then name the teacher, classroom, mark, and answer words that belong to this school setting.

\begin{grammarlens}[title={What you've built: check after speaking}]
A complete response has four connected turns and the setting words. One check is \textbf{\sk{परीक्षा} \sk{कदा} \sk{भविष्यति}? \sk{श्वः} \sk{परीक्षा} \sk{भविष्यति।} \sk{गृहकार्यं} \sk{करिष्यामि।} \sk{पाठं} \sk{पठिष्यामि।}} Then recall \textbf{\sk{विद्यालयः}, \sk{छात्रः}, \sk{अध्यापकः}, \sk{समयसारणी}, \sk{कक्षा}, \sk{अंकः}, \sk{उत्तरम्}}.
\end{grammarlens}

\begin{tcolorbox}[breakable,colback=teal!4,colframe=teal!35!black,title={Before you move on}]
What held the old school words together? (One timetable question, its answer, and two first-person plans.)
\end{tcolorbox}
\section[(independent school plan, second round)]{Independent school plan: round two}

\label{lesson:SA-C267-school-plan-retrieval}



\begin{quote}
Say “a hut,” then make one familiar near-term first-person plan. (\emph{Kuṭīram; use the taught śvaḥ/paraśvaḥ plan frame.})
\end{quote}

\subsection*{Your turn: a fresh prompt}
No Sanskrit model is shown before you answer.

Imagine that a teacher points to the timetable in a classroom. Ask when the exam will be and answer that it will be tomorrow. Say what homework you will do and what lesson you will study. Finish by recalling the school, student, teacher, mark, and answer words.

\begin{grammarlens}[title={What you've built: boundary}]
You can now complete one small A1 school-day plan with familiar language. This is \textbf{not} complete classroom discourse, school-system, grading, educational rights, or future-paradigm coverage, and it is not evidence of exam readiness, human validation, or learner success.
\end{grammarlens}

\begin{tcolorbox}[breakable,colback=teal!4,colframe=teal!35!black,title={Before you move on}]
Name the four jobs: read the timetable, ask and answer the exam time, state the homework plan, and state the study plan.
\end{tcolorbox}
